\subsection{3.3 智能体攀登}\label{agentic-climb}

智能体攀登训练 MAI-Thinking-1 解决需要与外部环境交互而非一次性文本应答的任务。在此设定下，模型必须拆解用户请求、选择工具或代码操作、观察结果，并在需要时跨多个步骤和轮次调整计划。训练信号将可验证奖励（如软件环境中测试通过、或数据库达到目标状态）与针对难以精确刻画的行为方面（如任务理解、有用性和轨迹质量）的 AI 反馈奖励相结合。

\pandocbounded{\includegraphics[keepaspectratio,alt={image}]{images/c4eb637ee18aed107b2611ffed3f863506e0355fc65175a9e77fc35a4b341ad2.jpg}}

图 18. RL 训练中的智能体循环与多步编排。

我们聚焦两个智能体领域：(i) 软件工程（SWE）（第 3.3.1 节），涉及基于真实代码库构建的可执行软件工程环境；(ii) 通用工具使用（第 3.3.2 节），涉及在多步任务中调用结构化工具。

在实践中，我们联合攀登由智能体任务和推理型 STEM 任务构成的混合任务集，其中包含我们的竞赛编程混合集。我们发现，纳入 STEM 任务有助于稳定 RL 攀登过程，并对多步软件工程和工具调用性能产生正向迁移。反之，智能体任务对 STEM 相关的单次通过（single-pass）性能既无正向也无负向迁移。与 STEM 攀登类似，我们在整个攀登过程中观察到模型行为不断演化，示例见附录 D.2。智能体场景与 STEM 场景的不同之处在于其使用多步 rollout 和容器化环境，具体如后续小节所述。

多步 RL 框架。我们的智能体多步 RL 与第 3.1 节中单步推理 RL 配方采用相同的核心目标，但将 rollout 从单个采样响应扩展为由策略步骤与环境步骤（观测）构成的轨迹。每个 RL 环境由一个任务规范、一个用于执行工具的沙盒执行环境（Sandbox Execution Environment，SEE）会话，以及一组可验证或经评判的奖励组成，用于在跟踪环境状态的同时评估任务完成情况。在每个策略步骤中，模型可以发出工具调用或产生最终答案。工具调用在 SEE 会话内执行，其输出会在下一个策略步骤之前追加到上下文中。随后对完整轨迹的正确性进行评分，并可选择性地给出 AI 反馈，之后在所有策略步骤产生的所有 token 上统一进行信用分配。

编排框架是一个 ReAct 式循环（Yao et al., 2023）：它解析模型的推理和动作，将工具调用分派给 SEE，把返回的观测追加到上下文，并将控制权交还给策略以进行下一步动作。之前所有步骤的 token 都保留在上下文中，并作为后续步骤的严格前缀。同一循环同时支持软件工程和通用工具使用两类工作负载。在 SWE 中，动作包括读写文件、运行 shell 命令，以及在针对问题定制的容器镜像内检查仓库状态。在通用工具使用中，动作调用由可变任务状态（如带种子数据的数据库）支撑的结构化工具，工具输出作为观测返回，供模型下一步动作使用。当模型不再发出工具调用，或 rollout 超过步骤、上下文长度或时间预算时，循环终止；随后 rollout 及其对应的 SEE 被送往评分器。评分器综合运用格式检查、基于规则的检查、可执行测试、可验证的状态比较和 AI 裁判。图 18 概括了我们的智能体 RL 框架。RL 环境来自两大领域：软件工程（第 3.3.1 节）和通用工具使用（第 3.3.2 节）。

沙盒执行环境。我们的沙盒执行环境（SEE）按需大规模提供隔离的容器环境。SEE 在 RL 训练期间用于为模型训练环境提供高度并行、低延迟的容器。它为每个智能体任务启动一个新容器，并在任务完成后销毁该容器。这保证了可复现性，支持安全探索，并防止任务之间的状态泄漏。

容器默认进行网络隔离，确保训练回合（episode）具有确定性，不受速率限制或瞬时故障等外部副作用的影响。当环境确实需要网络访问（例如安装软件包）时，流量通过缓存代理和域名白名单进行管控，以平衡环境保真度、训练可复现性与安全性的需求。

\subsection{3.3.1 软件工程}\label{software-engineering}

训练前沿模型执行 SWE 任务，需要模型在真实的开发环境中与真实代码库交互。大量高质量的编码环境对于训练强大的编码模型至关重要。为实现这一目标，我们从真实代码库的贡献中收集自然产生的数据。为整理我们的 RL 环境库，我们依赖以下基础组件：大规模代码仓库数据集（如公开的 GitHub 代码）、从真实世界开源开发者的 pull request 构建可执行软件工程问题的流水线，以及用于执行的沙盒容器服务 SEE。

SWE RL 环境、任务与工具。每个 SWE RL 问题被打包为自包含的容器镜像，其中包含检出在特定 commit、预装全部依赖的仓库、问题描述以及用于评分的单元测试。由于镜像捕获了仓库的全部依赖，环境具有确定性，且 rollout 开始时无需任何配置即可立即使用。在 rollout 过程中，模型通过工具调用与容器交互：读取和编辑文件、运行 shell 命令以及浏览仓库。当模型发出完成信号或达到轮次上限时，评分器在同一个容器内执行测试，并将结果与预期输出进行比较，产生可验证的奖励信号。

模型可以使用两种常用工具：Bash——一种执行 bash 命令并返回标准输出和错误流的工具；以及 String replace editor——一种让模型通过精确字符串替换来编辑文件的工具，无需为此使用不符合人体工学的 bash 工具（Anthropic, 2025）。完整的工具描述和参数 schema 见附录 E。

自研可扩展 SWE 环境构建。已有工作（Pan et al., 2024; Chen et al., 2026b; Fu et al., 2026）探索了将来自互联网的原始代码仓库转化为 RL 环境。我们以公开 GitHub 仓库的 issue 和 pull request（PR）作为原始数据来源，并受 SWE-ReBench（Badertdinov et al., 2026a,b）启发，开发了一条可扩展的流水线，用于摄入真实世界可验证的 SWE 问题。为构建高质量的 RL 环境，我们确保问题描述为智能体正确实现解决方案提供足够的信息，同时不过度限定任务。我们还确保评分器提供足够的覆盖范围，同时避免过于具体的评估标准。流水线由以下阶段组成，更多基础设施细节见附录 G：

• 公开 GitHub PR 与 issue 收集。流水线从 1.02 亿个公开 GitHub PR 开始。首先，我们过滤 PR，以获得适合构建环境的子集。要入选，PR 必须已合并到其仓库的主分支，且修改的文件少于 15 个。每个 PR 还必须同时包含代码变更和测试变更，以便评分器将测试变更用作隐藏测试，评估模型修复问题的能力。测试变更和代码变更根据补丁中的文件内容进行拆分。我们进一步基于 issue 关联关系过滤 PR，仅保留与 GitHub、Jira、Bugzilla、YouTrack、Phabricator、Launchpad 或 Linear 的 issue 相关联的 PR。经过此阶段，我们获得约 487 万个带有关联 issue 的 PR。

• 自动智能体环境构建。我们使用一个读取仓库状态并创建 Docker 文件以构建可执行容器镜像的 LLM 智能体，将选中的 PR 转换为可训练的环境。我们通过执行测试套件来验证每个镜像，并丢弃存在依赖或环境错误的构建。

• 参考评分信号提取。对每个候选问题，我们以 base commit 为基准分两个阶段执行仓库测试套件，得出参考测试结果：首先仅应用测试 diff（修复前，pre-fix），然后同时应用测试和代码 diff（修复后，post-fix）。由失败转通过（fail-to-pass，F2P）的测试构成问题解决信号，要求模型补丁使这些测试反转通过，才能获得解决问题的分数。在两个阶段均保持通过的测试，即持续通过（pass-to-pass，P2P），构成回归信号，确保生成的补丁不会破坏现有功能。没有存活 F2P 测试的问题被丢弃。

• 环境与评分器验证。在获得可执行环境、测试和代码变更、测试执行脚本以及预期的修复前后结果后，我们使用与 RL 训练期间相同的 SEE 训练基础设施对环境进行重新验证。这是对环境构建沙盒与训练沙盒之间潜在差异的最终合理性检查。我们验证空代码补丁无法通过评分器，而标准解法补丁能够通过，且各自经过多次试验。尽管上游补丁已知是正确的，但重新验证仍可能因集群差异（如网络出口策略、CPU/内存限制、执行超时）而失败。我们还过滤在重复执行中表现出非确定性测试行为的环境，以降低奖励噪声。

• 质量过滤与 SWE 环境改写。仅靠正确执行和严密的测试还不够，因为许多环境仍包含低质量或定义不充分的问题描述。例如，我们经常看到简短含糊的 bug 描述、不完整的规格说明，或类似“修复问题”这种只在隐藏测试中隐式编码了预期行为的描述。如果问题描述与评分器最终运行的测试所期望的内容之间存在较大差距，模型就只能猜测测试需要什么。

为提升环境质量，我们在同一环境中部署一个智能体，检查问题描述、仓库和测试，并从规格清晰度、测试质量、泄漏风险和可行性等方面对任务进行评分。对低质量环境，智能体会改写问题描述，使其更好地与测试要求对齐，同时避免不必要的信息泄漏或过度限定。

在最初的 487 万个候选问题中，208 万个（42.8\%）成功通过自动智能体环境构建，745,452 个（15.3\%）通过参考评分信号提取，265,617 个（5.5\%）在 94,044 个独立仓库中通过了环境与评分器验证。

合成数据。虽然我们的环境构建流水线生成的许多环境未通过验证，但其中很大一部分仍然是有效的可执行环境。在大多数情况下，失败源于低质量的问题描述或不足的测试覆盖，而非无效的执行环境。为复用这些可执行环境，我们借鉴 BugPilot（Sonwane et al., 2025）、SWE-Smith（Yang et al., 2026）和 SWE-Mirror（Wang et al., 2025a）的方法，从与成功构建但未通过质量检查的环境相关联的 PR 中生成新的合成问题及相应测试。

防止 RL 训练中的奖励操纵。尽管 SWE 环境使用可验证、可执行的测试用例进行评分，它们仍然容易受到奖励操纵（reward hacking）的影响。我们使用 LLM 监控器审查 rollout，并对标记的示例进行人工审查。这识别出三种主要的奖励操纵类型：

• 互联网搜索。由于我们的 RL 环境来自 GitHub 上的公开仓库，PR（以及其中的标准解法）往往可以通过搜索获取。我们对沙盒执行环境实施网络访问控制来防止这种情况：当问题自包含时禁用互联网访问，否则仅允许问题所需的最低限度访问。

• 本地 git 历史搜索。Git 是软件仓库不可或缺的组成部分。有时智能体会搜索所有 git commit 和日志，试图找到隐藏在本地 git 数据库中的解答 commit。\footnote{例如，SWE-Bench 曾受此问题影响：https://github.com/SWE-bench/SWE-bench/issues/465。此后该问题已修复。}删除所有 git commit 也不可行，因为 git 是一项有价值的技能，并能提供解决问题所需的正当信息。我们通过清除问题 base commit 之后的所有 commit、引用和分支来彻底净化环境，创建仓库“穿越”回原始状态的版本。

• 篡改测试。遵循 SWE-Bench（Jimenez et al., 2024）使用的评估协议，我们在评分前重置所有被智能体修改的测试文件，防止智能体篡改测试用例使其通过评分。我们还在推理期间对智能体隐藏测试变更，仅在评分时应用，以进一步减少作弊机会。然而，仍然存在作弊途径，例如通过 monkey-patching 测试框架或修改等价性行为，这些难以通过简单地重置测试文件或引入隐藏测试来预防。我们使用 LLM 监控智能体行为，并持续加强测试文件检测、测试文件重置及其他反篡改启发式策略。

\subsection{3.3.2 通用工具使用}\label{general-tool-use}

与 SWE RL 环境相比，通用工具使用 RL 环境在可用工具和应用领域两方面都表现出显著更大的多样性。它们涵盖广泛的通用工具调用任务（例如库存管理、调度平台、报告生成、客户支持及其他面向企业的场景），强调与外部系统进行可靠、有状态的交互。

工具使用 RL 环境。每个工具使用 RL 问题被实例化为一个交互式、有状态的环境，由模拟真实 API 或 MCP（Model Context Protocol，2026）行为的模拟后端支撑。具体而言，每个问题由一条查询、一组带 schema\footnote{我们使用 OpenAI 的 function-calling schema：https://developers.openai.com/api/docs/guides/function-calling。}的可用工具、一个初始环境状态和一个评分器组成。与通常只涉及少量工具的 SWE 设置不同，工具使用环境以远为庞大的工具集模拟丰富的真实服务交互，单个环境中的工具数量往往超过 50 个。这一设定训练模型高效地选择合适的工具，并改善跨多样化工具使用场景的泛化能力。我们同时使用人工整理的环境和由我们的工具调用任务生成框架（下文介绍）生成的合成环境进行训练。

合成环境。为扩大训练数据的规模和多样性，我们合成生成自包含的封闭世界环境，包括带种子数据的数据库、工具定义和可验证的任务。我们的流水线只需对目标环境进行简单英文描述，即可端到端生成完整的工具使用环境。该方法延续了自动生成工具使用环境的已有工作（Sullivan et al., 2025; Zhang et al., 2026; Wang et al., 2026; Mika, 2026），特别是 FunReason-MT 流水线（Xu et al., 2025b; Hao et al., 2025）。总体而言，我们的目标是复现常见的企业和消费者工具使用场景，例如旅行预订和库存管理。

我们的流水线包含三个主要阶段：(i) 环境引导——生成工具描述、实现函数，并用相关实体为数据库填充种子数据；(ii) 任务创建——采样可能的工具调用轨迹、创建与交互链相关的实体并构建用户请求；(iii) 验证与优化——执行上一阶段生成的动作，并移除过于相似的任务。每个阶段都涉及多次 LLM 调用，用于实体生成和函数实现等任务。我们还采用迭代的批判与改进循环来优化每个阶段的输出（Dhuliawala et al., 2023; Madaan et al., 2023）。

我们生成环境特定的用户人设（personas）（Ge et al., 2024），以进一步丰富任务的多样性。为缓解过度积极的工具调用行为，我们还用包含工具描述但不需要使用工具的任务来扩充数据。总体而言，我们的流水线合成了超过 150 个环境和 130,000 个任务。

奖励设计。训练期间，我们混合使用环境专属评分器和跨环境评分器。环境专属评分器根据最终环境状态、工具使用模式和最终答案分配奖励。对合成环境，我们使用 LLM 裁判将任务拆分为多个子任务并分别评分。我们还采用鼓励高效工具使用的跨环境评分器，包括尽可能并行调用工具、避免重复调用，以及以有效的参数类型和参数正确调用工具。
